\section{Application}\label{sec:application}\label{sec:simulation}

To assess how reliably finite samples of financial news recover the positive association
predicted by the model, we develop a multi-agent simulation of the economy in
Section~\ref{sec:model}, whose population coefficients are known.
We examine the effects of article coverage and changes in reporting on these
estimates, then apply the estimator to corporate news in
Section~\ref{sec:data}.

\subsection{Experimental Design}\label{sec:simulation-design}

The experiment follows competitive adjustment after small exogenous free
cash flow shocks,
keeping operating movements close to the local setting of the directional
prediction.
Sampling articles along these paths reveals how well the estimator detects
small reallocations at the coverage observed in the corporate data.

The competitive economy has 100 firms with canonical interior baseline
shares $p_i=i/101$, total
capacity one, two activities and two opportunities.
Outside appeal and baseline opportunity levels equal one, with cost curvature
$\nu=1/(4D^2)$, where $D=51$, linear costs supporting the interior origin,
and adjustment weight $\kappa=d_\perp=\nu+1/D^2$.

The 60 monthly ticks begin from stationary conditions, which we label 2018 for
comparison with the corporate sample, followed by independent
permanent innovations of $\pm0.001$ at ticks 18, 30, 42 and 54 that raise one
opportunity and lower the other.
We evaluate all sixteen sign histories and assign each equal weight.
At monthly discounting $\delta=0.95^{1/12}$, we compute return responses using
the infinite permanent-response resolvent in \eqref{eq:same-path-value}, while the
operating engine solves every contemporaneous joint Nash stage at the actual
positive opportunity levels.

In the first responding month after the first innovation, absolute changes
in normalized operating shares range from
\ppnum[4]{ir-redraw-signal-first-response-min-pct} to
\ppnum[4]{ir-redraw-signal-first-response-max-pct}
percentage points across firms and histories, with the largest absolute departure
from an origin share over the complete paths reaching
\ppnum[3]{ir-redraw-signal-maximum-origin-change-pct} percentage points.
These small movements specify the local signal tested in this economy.

The baseline article budget uses the firm-year article counts and frozen
empirical monthly date distributions of the primary corporate sample
described in Section~\ref{sec:data}.
The next budgets multiply every count by four and sixteen.
For each history, \ppnum[0]{ir-validation-redraws} independent article panels
use nested count increments
across budgets and shared sampled monthly counts across reporting regimes.
The same dated counts determine annual article laws and article-timing
weights, and each sampled origin law is reused wherever it serves as a
reference.
Because multinomial month and atom counts can leave unequal masses on antipodal
atoms, we compute transport exactly on the axis support, which retains that
asymmetry.
Reusing each solved operating path across redraws isolates article
measurement variation around the same economic response.

For each history $h$, we distinguish approximation of the operating response
from its measurement through articles.
Comparing the affine-share coefficient with the finite-capacity
common-reporting coefficient $\beta_h^0$ assesses the local economic
approximation.
The population reporting-law coefficient $\beta_h^\lambda$ incorporates
changing attention, while the sampled estimate
$\widehat\beta_{h,r,s}^\lambda$ adds finite-article variation at budget $s$.
Table~\ref{tab:directional-validation-controls} reports absolute error
against each history's population-law target and the frequency of numerically
resolved positive estimates, averaging histories equally.
Because every common-reporting population coefficient is positive, its
positive-estimate frequency also measures correct-sign recovery.
The fixed draw count targets a stratified aggregate rate standard error below
one percentage point.
Appendix~\ref{app:recovery-diagnostics} specifies the arithmetic sign allowance
and retains the full directional and magnitude diagnostics.

\subsection{Estimator Recovery and Reporting Effects}\label{sec:simulation-recovery}

% AUTO-GENERATED from saved Paper 6 results; do not edit.
\begin{table}[htbp]
\centering
\caption{Mechanism and measurement controls. Article-count multipliers scale observed coverage. Mean target is the arithmetic mean of history-specific population coefficients. Positive reports numerically resolved positive estimates as percentages (stratified Monte Carlo SE in percentage points). Absolute RMSE uses each history's matching target in coefficient units. Stationary reporting is the zero-alignment measurement control; reporting only holds capacities fixed and uses reporting amplitude 0.005. Table~\ref{tab:directional-validation-details} provides full diagnostics and definitions. Each percentage uses \ppnum[0]{ir-validation-common-1-positive-denominator} attempted redraws.}
\label{tab:directional-validation-controls}
\begin{tabular}{@{}l@{\enspace}r@{\enspace}r@{\enspace}r@{\enspace}r@{}}
\toprule
Design & \shortstack{Article-count\\multiplier} & Mean target & \shortstack{Positive estimates\\(\%)} & Absolute RMSE \\
\midrule
Common reporting & $1\times$ & \ppnum[3]{ir-validation-common-1-mean-target} & \ppnum[1]{ir-validation-common-1-positive-pct}(\ppnum[2]{ir-validation-common-1-positive-se-pct}) & \ppnum[3]{ir-validation-common-1-rmse} \\
Common reporting & $4\times$ & \ppnum[3]{ir-validation-common-4-mean-target} & \ppnum[1]{ir-validation-common-4-positive-pct}(\ppnum[2]{ir-validation-common-4-positive-se-pct}) & \ppnum[3]{ir-validation-common-4-rmse} \\
Common reporting & $16\times$ & \ppnum[3]{ir-validation-common-16-mean-target} & \ppnum[1]{ir-validation-common-16-positive-pct}(\ppnum[2]{ir-validation-common-16-positive-se-pct}) & \ppnum[3]{ir-validation-common-16-rmse} \\
Stationary reporting & $1\times$ & \ppnum[3]{ir-validation-stationary-1-mean-target} & \ppnum[1]{ir-validation-stationary-1-positive-pct}(\ppnum[2]{ir-validation-stationary-1-positive-se-pct}) & \ppnum[3]{ir-validation-stationary-1-rmse} \\
Reporting only & $1\times$ & \ppnum[3]{ir-validation-reporting-only-mean-target} & \ppnum[1]{ir-validation-reporting-only-positive-pct}(\ppnum[2]{ir-validation-reporting-only-positive-se-pct}) & \ppnum[3]{ir-validation-reporting-only-rmse} \\
\bottomrule
\end{tabular}
\end{table}


Under common reporting, correct-sign recovery rises from
\ppnum[1]{ir-validation-common-1-correct-pct}\% at observed coverage to
\ppnum[1]{ir-validation-common-4-correct-pct}\% at fourfold coverage and
\ppnum[1]{ir-validation-common-16-correct-pct}\% at sixteenfold coverage.
Even the largest article samples therefore often fail to recover a
numerically resolved positive sign, despite a normalized difference of only
\ppnume[2]{ir-finite-affine} between the finite-capacity and affine-share
population coefficients.

The small operating-share changes shift population article weights only
slightly, so finite-article variation can obscure the corresponding movement
toward reference profiles.
Coefficient error also falls with article coverage, but remains large
relative to the population target.
Normalized RMSE under common reporting is
\ppnum[2]{ir-validation-common-1-normalized-rmse},
\ppnum[2]{ir-validation-common-4-normalized-rmse} and
\ppnum[2]{ir-validation-common-16-normalized-rmse} across the same ladder.
Even at sixteenfold coverage, error remains several times the root-mean-square
magnitude of the population coefficients.

Reporting intensity can also respond to the opportunity itself.
We set $\lambda_{\pm,t}=\exp(\pm0.005u_t)$, where $u_t$ is cumulative
news divided by $0.004$, its maximum possible magnitude, so attention shifts
toward the activity whose contribution-margin opportunity has increased.
At sixteenfold coverage, correct-sign recovery relative to the operating-law
target rises to
\ppnum[1]{ir-validation-changing-16-correct-pct}\%
(Appendix~\ref{app:recovery-diagnostics}).
At fixed capacities, reporting changes alone generate a positive population
coefficient in all
\ppnum[0]{ir-reporting-only-positive-histories} histories.
The reporting-only row in Table~\ref{tab:directional-validation-controls}
measures this reported-law target, while its operating-law target is zero.
The counterfactual retains return responses from the modeled operating
path and changes only the observation law, without solving a separate
fixed-capacity economy.

Holding the reported population laws stationary under common reporting
isolates article noise around a zero-alignment target.
We retain the dated financial process, so population approach is exactly zero
while the financial denominator remains positive.
Finite articles nevertheless yield coefficients of both signs, with positive
frequencies near one half across the coverage ladder.
Because a zero target supplies neither a correct positive sign nor a scale
for normalized RMSE, we assess this measurement control through absolute
errors and sign frequencies.
The control holds reporting laws fixed rather than solving a new competitive
equilibrium.

\subsection{Measurement Uncertainty and Inference}\label{sec:simulation-uncertainty}

The spread of finite-article estimates provides a benchmark for assessing
whether redrawing articles from an observed archive represents that
measurement uncertainty.
We generate independent archives from the common-reporting, stationary and
reporting-only laws, retaining 100 firms and observed article coverage, then
apply the corporate article-redraw procedure within each archive.
Across \ppnum[0]{ir-redraw-outer-replications} archives per history, each
with \ppnum[0]{ir-redraw-inner-redraws} inner redraws, the broad redraw
ranges give near-nominal containment of the population coefficients.
Containment and detection answer different questions.
Under the positive-target common-reporting and reporting-only laws, the fifth
percentile exceeds zero in
\ppnum[1]{ir-redraw-mechanism-positive-q05-pct}\% and
\ppnum[1]{ir-redraw-reporting-positive-q05-pct}\% of attempts, respectively.
Thus, intervals that broadly capture measurement uncertainty rarely distinguish
these small positive population coefficients from zero.
Under stationary reporting, the target is zero, so the corresponding
\ppnum[1]{ir-redraw-stationary-positive-q05-pct}\% is instead a null-control
rejection frequency.

Uncertainty across independently sampled firms concerns a different target
from article selection around fixed operating paths.
For each target coefficient, we generate
\ppnum[0]{ir-inference-n100-replications} independent panels, each consisting
of iid firm histories.
Increasing the number of firms from 16 to 100 raises coverage of the
candidate one-sided 95\% lower bound from
\ppnum[1]{ir-inference-audit-n16-coverage-pct}\% to
\ppnum[1]{ir-inference-n100-coverage-pct}\%.
Zero-null rejection falls from
\ppnum[1]{ir-inference-audit-n16-null-rejection-pct}\% to
\ppnum[1]{ir-inference-n100-null-rejection-pct}\%
(Monte Carlo SE \ppnum[2]{ir-inference-n100-null-rejection-se-pct}
percentage points at 100 firms).
These generated-history panels preserve dependence across years and
comparison roles under the specified iid sampling law, although they are not
jointly solved competitive economies.
Appendix~\ref{app:annual-gmm} gives the sampling designs, target containment
and Monte Carlo errors.
